Beat policing---i.e., assigning the same officers to smaller, well-defined territorial units rather than rotating them over larger ones---has become a widespread strategy to address crime and maintain public safety, yet rigorous causal evidence on its effectiveness remains scarce. Using administrative and survey data from Chile, we show that \emph{Plan Cuadrante}, a reform that introduced beat policing in urban municipalities, substantially reduced victimization and crime perception, while increasing trust in the police. Although its implementation required increasing police resources in some municipalities, we show that resource increases alone cannot account for the observed improvements. The beat policing elements of the reform---permanent assignment to quadrants enabling officers to build local knowledge---seem to be key drivers of its effectiveness.

The reform not only reduced crime incidence but also insecurity perceptions. As a result, households in adopting municipalities reduced their investment in private security measures. Crime rates and perceived victimization risk do not always move together \citep{Ajzenman_Pato}, yet \emph{Plan Cuadrante} reduced both simultaneously. This is an important result as misalignments between actual and perceived risk can lead households to overinvest in private security and distort public expenditure priorities \citep{Minali, Ajzenman_Pato}.

The reform also increased trust in the police---an outcome that has proven particularly difficult to achieve in low- and middle-income countries. Community policing, long regarded as the primary strategy for building police legitimacy, has shown limited effects on both crime and trust in such settings \citep{blair2021, Tobon}. Beat-policing takes a different approach. Rather than relying on community co-production of security, it builds proximity and trust organically through sustained, stable territorial engagement. Our results suggest this organizational logic can be effective even outside developed-country contexts.

Taken together, these results contribute to a broader body of evidence that highlights that the spatial organization of police resources matters. By assigning officers permanently to smaller territorial units, beat policing increases visibility, deepens local knowledge, and improves the capacity of officers to prevent crime and respond to it effectively. This pattern echoes findings from health and education showing that aligning frontline service providers more closely with the communities they serve enhances outcomes \citep{besley2003centralized, Bardhan, Bloom_2015, Basurto}.

While \emph{Plan Cuadrante} was implemented under specific conditions, those conditions help clarify when beat policing is likely to succeed. Two institutional features appear particularly important. First, the absence of significant organized crime during the study period allowed officers to operate repeatedly in their assigned areas without meaningful risk of intimidation or capture, enabling local knowledge and community ties to take root. Second, the relatively high level of professional training in the Chilean police provided a foundation for effectively implementing the reform. Neither of these features is idiosyncratic. They describe an institutional environment that is present, or could be developed, in many settings. Together, they suggest that the gains from beat policing are most likely to materialize when supported by a capable police force operating in a security environment where persistent local engagement is feasible.